\section{Data and the Spawn Clock}
\label{sec:data}

\tab{tab_data}
\tab{tab_rules}

\paragraph*{Games}
The core set is every game of the 2025 Mid-Season Invitational and the 2025 World Championship, \V{nGames} games in \V{nSeries} series between \V{nTeams} teams (Table~\ref{tab:data}).
Both events were played on patches 25.13 and 25.20 under the same objective rules, so one spawn clock covers all games.
For each game we hold three sources: the broadcast video, Riot's match-v5 timeline~\cite{riotapi} with exact positions of all ten players every 60\,s and millisecond-stamped kills and objectives, and Leaguepedia~\cite{leaguepedia} metadata (roles, draft, series).
An archive of 2,678 professional games with v5 timelines (2023--2026) is used only to fit outcome models: the win surface, event values and the navigation graph.
It never sees a setup choice.

\paragraph*{The spawn clock}
Table~\ref{tab:rules} lists the rules.
We infer every spawn time from the rules and the kill events and check that no kill precedes its inferred spawn; none does across the \V{nSpawns} spawns.
Anchoring at the spawn rather than at the kill is the central design choice of this paper.
The spawn time is fixed by rule and known to both teams, so a window anchored there does not depend on what the teams did; a window anchored at the kill selects on who took the objective and when.
After the spawn we keep a team-snapshot only while the objective is still alive (a risk set), because a snapshot taken after the kill describes the aftermath rather than the setup.

\paragraph*{Units}
A \emph{team-setup} is one team at one spawn, \V{nSetups} in total.
We observe it on a grid of snapshots every 15\,s from $s{-}120$ to $s{+}180$\,s, where $s$ is the spawn time.
Two layers share these units.
The \textbf{V layer} uses 2-fps minimap tracks from the broadcast (Section~\ref{sec:measurement}) and covers every spawn and every snapshot.
The \textbf{R layer} uses only Riot's exact 60-s frames and covers the fixed-time first spawns (dragon, grubs, herald, Atakhan, baron) at $s{-}60$, $s$ and $s{+}60$.
The R layer has no measurement error and coarse time; the V layer has fine time and validated measurement error.
Findings that matter are checked on both.

\paragraph*{Gold state}
A team is \emph{trailing}, \emph{even} or \emph{leading} by its gold difference at $s{-}120$\,s divided by the archive standard deviation of the gold difference at that game minute, with cut points $\pm 0.5$.
The state is fixed before the setup begins, so it cannot respond to the setup.

\paragraph*{Objective classes}
Mixture weights in the discovery layer are stratified by three classes that differ in timing and map: early river objectives (dragons, grubs, herald), Atakhan, and late objectives (baron, elder).
In the descriptive analysis respawned dragons and barons form their own groups because their spawn time depends on play.
